% !TeX root = ../main.tex

\chapter{数值实验和结果分析}

为了全面验证本文提出的基于线图映射的非加性最短路径算法（CWSP）的逻辑正确性，以及面向 GPU 架构的并行优化方案的加速效能，本章将开展多组定性与定量的对比实验。通过定制化微观路网场景与大规模模拟网络拓扑，分别从算法的求解可行性与超大规模并发吞吐两个维度进行深度剖析。

\section{实验准备}
% 介绍你的电脑配置、代码运行的环境，以及使用的数据长什么样。

\subsection{软硬件运行环境}
% 说明公平对比的基准：串行算法运行的 CPU 环境与语言（如 Python/NetworkX），以及 GPU 算法运行的算力平台与底层框架（如 CUDA/cuGraph/CuPy）。

\subsection{实验数据集说明}
为了全面且客观地评估本文算法的适用性与扩展性，本章实验同时采用了以下两类数据集进行交叉验证：
\begin{itemize}
    \item \textbf{自建随机微观网格数据集：} 主要用于定性验证与算法可行性检查。通过编写脚本随机生成包含特定约束（如“禁止特定方向左转”、“某连贯边存在高昂惩罚”）的 $N \times N$ 定制网格拓扑。这类数据的绝对路径与边约束完全已知，便于人工核对 Dijkstra 算法与 CWSP 算法的输出路径差异。
    \item \textbf{真实世界大规模路网数据集（Real-world Road Networks）：} 为了与顶尖文献\cite{durand2024generalized}中的压力测试条件对齐，并在绝对算力层级上激发 GPU 的并行潜能，本章引入了基于真实地理状态提取的大规模公开路网基准数据集（如源自 OpenStreetMap 提取的城市道路图或 第9届 DIMACS 最短路径挑战集提供的真实路网资料，规模从数万至数百万节点不等）。原图附带了真实路口的夹角转弯、通行级别等多元物理维度，能够充当最苛刻的非加性网络高频测试床。
\end{itemize}


\section{实验设计与评价指标}
% 告知读者对比的维度与指标设定。

\subsection{对比算法基准设置}
% 说明你的三位“参赛选手”：
% 1. 标准 Dijkstra 算法（不考虑前驱约束的基准）
% 2. 基于串行 CPU 的 CWSP 算法（本文理论基础，处理非加性约束但速度慢）
% 3. 基于 GPU 并行优化的 CWSP 算法（本文核心创新，即包含并行线图构造与 Delta-Stepping 的版本）

\subsection{性能评价指标}
% 提出评判标准：1. 路径可行性（是否违规/是否真实最优）；2. 绝对计算耗时 (ms)；3. GPU 相对 CPU 的加速比。


\section{基于线图的串行算法与传统最短路径算法对比}

\subsection{路径可行性验证场景构造}
为了直观展示传统最短路径算法在非加性网络（如带转弯惩罚的路网）中的局限性，本节基于前期提取的迈阿密国际机场（MIA）滑行道真实拓扑，构造了极具代表性的微观局部陷阱。在航空地质调度中，飞机如果在高滑行速度下强行进行锐角转弯存在极大的事故风险。因此实验配置了严苛的转向约束条件：当两条相邻连通边构成的转弯夹角过大（如超过一定阈值构成了物理上的锐角转弯），将触发高达 $M=5000$ 米甚至正无穷大的硬性惩罚代价。

典型地，在迈阿密机场网络中，通过程序挖掘定位了起点节点 $v_s$（ID：2617）与终点节点 $v_d$（ID：2988）之间的一个死端约束枢纽 $v_m$（ID：1483）。在纯几何距离视角下，最短的连通方案为 $v_s \rightarrow v_m \rightarrow v_d$，其空间物理直线距离仅有 $48.43$ 米。然而，这一几何短路正是由两条近乎首尾折返的边相接而成，该走势在该路口构成了一个非法的“掉头式转角”。

\subsection{解空间与路径代价对比分析}
为验证上述理论分析，将经典的 Dijkstra 最短路径算法与本文探讨的基于线图拓扑映射的串行 CWSP 算法（使用 Python 原型实现进行逻辑基准标定）置于该机场网络的多个随机测试用例上测算，结果汇总于表~\ref{tab:baseline_compare} 中。值得注意的是，实验同时考察了传统 Dijkstra 在“无视转弯约束（仅基于距离）”与“叠加物理可行性约束（若触发转弯限制则立刻视为无效状态置为 $+\infty$）”两种情境下的表现。

\begin{table}[htbp]
    \centering
    \caption{传统 Dijkstra 与线图 CWSP 在多组真实机场滑行道测试对中的路径代价值对比（单位：米）}
    \label{tab:baseline_compare}
    \begin{tabular}{cccc}
        \toprule
        \textbf{测试点对编号} & \textbf{\makecell{传统 Dijkstra 算法\\(无视约束，纯地理短路)}} & \textbf{\makecell{传统 Dijkstra 算法\\(考虑可行性约束限制)}} & \textbf{\makecell{基于线图的非加性 CWSP\\(严格遵循约束的合法短路)}} \\
        \midrule
        1  & 56.6    & $+\infty$ & 472.7  \\
        2  & 46.8    & $+\infty$ & 349.8  \\
        3  & 1624.1  & 1624.1    & 1624.1 \\
        4  & 844.1   & $+\infty$ & 920.2  \\
        5  & 766.5   & 766.5     & 766.5  \\
        12 & 27.7    & $+\infty$ & 2396.4 \\
        约束陷阱点对 & 48.4    & $+\infty$ & 548.8  \\
        \bottomrule
    \end{tabular}
\end{table}

综合表~\ref{tab:baseline_compare} 中提炼出的一系列极端情况与顺畅情况指标，可以给出支撑本课题理论可行性的三大核心定性判断：
\begin{enumerate}
    \item \textbf{传统贪心策略天然无法应对非加性约束：} 在带有多边连贯约束的情境下，如图中测试点对1、2、12 及特地挖掘出的陷阱点对所示，Dijkstra 算法受限于马尔可夫无后效性质，仅基于节点累积权重进行最短前驱选取。其视野局限于相距仅 $27.7$ 米或 $48.4$ 米的所谓“最优”近点。然而，其推导出的路径在现实中不可避免地会碰撞极为高昂的夹角惩罚（导致可行解变为 $+\infty$ 的完全失效）。
    \item \textbf{自然畅通路段等效性与鲁棒性：} 观察无锐角挑战的自然开阔点对（第3和第5组测试），原本不存在违规约束带来的阻碍。此时传统策略与线图变换后的 CWSP 解题结果分毫不差。这一组对照表明算法具备向下兼容属性，使用线图拓扑并未改变图中正常的流通距离基准。
    \item \textbf{线图模型能跨越障碍锁定全局严格最优：} 在陷阱点对（2617$\rightarrow$2988）的求解过程中，线图 CWSP 算法完美发挥了其“将 1-加性转弯惩罚静态化为 $L(G)$ 顶点权”的特性，并由此成功实现全局统筹规划。当探知枢纽处存在 $5000$ 级别违规惩罚时，算法驱动路径避开此瓶颈，并在外围一系列滑行道中划取了一个全角度合法的超级环形闭合圈。算法虽付出了 $548.8$ 米（原物理最短路 $48.4$ 米代价的近$11$倍）的极度周转长度，却在数学与物理上成功避免了非法后果的发生，顺利摘取到具有高度安全保障条件下的全局合法非加性最短解。这一机制也正是线图拓扑变换降维打击的核心所在。
\end{enumerate}


\section{GPU并行算法与串行算法的性能对比与加速比分析}
% 核心实验二：证明“线图导致规模爆炸时，GPU 并行策略的降维加速能力”。

\subsection{不同路网规模下的计算总耗时统计}
% 画折线图。横坐标是原图规模，纵坐标是对数尺度的耗时。展示规模爆炸时 CPU 的指数级卡顿和 GPU 的强劲吞吐。

\subsection{线图拓扑转换阶段的加速比分析}
% 呼应第三章 3.3.1 节。用柱状图对比“CPU 动态建图” vs “GPU 基于前缀和与 CSR 格式的并发写入”的时间差距。

\subsection{寻路松弛阶段的加速比分析}
% 呼应第三章 3.3.3 节。对比“原版基于优先队列的寻路”和“基于 GPU Delta-Stepping 算法”的性能差异。


\section{本章小结}
% 总结本实验得出的结论。

